\documentclass[10pt,a4paper]{amsart}
\usepackage[margin=19mm]{geometry}
\usepackage{amsmath,amssymb,mathtools}
\usepackage[scr=rsfs,cal=euler]{mathalfa}
\usepackage{xfrac}
\usepackage[unicode,hidelinks,bookmarks=false]{hyperref}
\newcommand{\ie}{i.e.\ }
\newcommand{\cf}{cf.\ }
\newcommand{\resp}{respectively\ }
\newcommand{\Subsection}{\subsection}
\providecommand{\nobreakdash}{\nobreak\hbox{-}}
\setlength{\emergencystretch}{2em}
\multlinegap0pt
\usepackage{bm}
\usepackage{bbm}
\usepackage{dsfont}
\usepackage{xspace}
\usepackage{cancel}
\usepackage{mathtools}
\usepackage{amsthm}
\makeatletter
\@ifundefined{theorem}{\newtheorem{theorem}{Theorem}}{}
\makeatother
\title[Framingtopes]{Framingtopes}
\author{Sergio Alejandro Fernandez de soto Guerrero, Cesar Ceballos, Matias von Bell}
\date{Source: arXiv:2609.20656v1 (CC BY 4.0)}
\begin{document}
\maketitle

\noindent\emph{Source and licence.} Framingtopes. Version arXiv:2609.20656v1 (CC BY 4.0), distributed under the Creative Commons Attribution 4.0 International licence, as recorded in the arXiv licence metadata for this exact version and shown on the abstract page https://arxiv.org/abs/2609.20656v1. Full rights notice: licenses/arxiv-2609.20656-CC-BY-4.0.txt.\par\smallskip

\noindent\textit{Editorial scope.} Counted unit: the Theorem whose source label is \texttt{thm\_Cayley\_tropicalArrangement} in the article (source0.tex lines 2103--2109, source label \texttt{thm\_Cayley\_tropicalArrangement}), delivered together with the article's own complete original proof of it (source0.tex lines 2111--2197, one \texttt{proof} environment of 87 lines). No mathematics is written, repaired or re-derived: statement and proof are the article's own text line for line. Required context retained uncounted: none. Every definition, display and label of those lines is retained unchanged. Omitted from this package and not redistributed: 1--2102, 2110--2110, 2198--3204. Nothing inside a retained excerpt is omitted or altered: every retained body is delivered exactly as the article's lines are written, so no omission receipt of any kind accompanies this package. Citations: the retained excerpts carry no citation command at all, so no bibliography entry of the article is transcribed and no reference is redistributed. Reference adaptation: none was needed and none was made; every reference inside the delivered unit resolves inside it, so no unresolved reference or citation is produced. Printed slips kept literal: no printed word, symbol, hypothesis or number of the counted statement or of its proof was altered. Layout and class: the article's own class is \texttt{amsart} with packages including array, amssymb, amsxtra, latexsym, graphicx, psfrag, epsfig, ifthen, mathtools, mathrsfs, xy, bbm, bm, ulem; the wrapper is the lane's pinned \texttt{amsart} editorial wrapper at 10pt on A4, which restates the theorem-like environments the retained text needs and adds the article's own macro definitions that the retained text needs (none), with extra wrapper packages (bm, bbm, dsfont, xspace, cancel, mathtools). The delivered numbering is the unit's own. Assets: no retained span contains a figure environment, a graphics-inclusion command, a PDF-page inclusion, a table or a TikZ picture; no image file is copied into this package, the package declares no asset dependency and no image file is redistributed with it. Licence: version arXiv:2609.20656v1 of this article is distributed under the Creative Commons Attribution 4.0 International licence, as recorded in the arXiv licence metadata for this exact version and shown on the abstract page https://arxiv.org/abs/2609.20656v1. A full rights notice is delivered at licenses/arxiv-2609.20656-CC-BY-4.0.txt. Characters: every retained excerpt is pure ASCII, so the delivered engine, XeLaTeX with its default Latin Modern text font, reports no missing character.
\begin{theorem}\label{thm_Cayley_tropicalArrangement}
Let $\mathcal{S}^h$ be the regular subdivision of $C(P_1,\ldots,P_k)$ induced by a height function $h$, and let $\mathscr{T}^h$ be the tropical cell decomposition of the corresponding arrangement of tropical hypersurfaces. Then the map
\[
Q\longmapsto\psi(Q)
\]
is an order-reversing bijection between the faces of $\mathcal{S}^h$ containing at least one vertex from each copy $\{e_i\}\times P_i$ and the cells of $\mathscr{T}^h$.
\end{theorem}

\begin{proof}
The faces of $\mathcal{S}^h$ are the projections of the upper faces of the lifted Cayley embedding of $C(P_1,\dots,P_k)$. Equivalently, they are the faces minimizing linear functionals of the form
\[
(-1,\mathbf{y},\mathbf{x}),
\]
where $\mathbf{y}=(y_1,\dots,y_k)\in\mathbb{R}^k$ and $\mathbf{x}\in\mathbb{R}^d$. Thus, a face $Q$ consists of all vertices $\{e_i\}\times R$, with $R\in\operatorname{Ver}(P_i)$, for which
\begin{equation}\label{eq_minimizedfaces}
-h(\{e_i\}\times R)+y_i+\langle R,\mathbf{x}\rangle
\end{equation}
is minimized over all $R\in\operatorname{Ver}(P_i)$ and $i\in[k]$.

Suppose that $Q$ contains at least one vertex from each copy $\{e_i\}\times P_i$, and write
\[
Q=
\operatorname{conv}\bigl((\{e_1\}\times Q_1)\cup\cdots\cup(\{e_k\}\times Q_k)\bigr).
\]
Fix $i\in[k]$. Since the term $y_i$ in~\eqref{eq_minimizedfaces} is constant on the copy $\{e_i\}\times P_i$, the face $Q_i$ is precisely the face of $P_i$ minimizing the linear functional
\[
-h(\{e_i\}\times R)+\langle R,\mathbf{x}\rangle,
\qquad
R\in\operatorname{Ver}(P_i).
\]
Hence,
\[
\mathbf{x}\in\tilde Q_i=\psi_i(Q_i),
\]
so each $\tilde Q_i$ is non-empty. Therefore,
\[
\mathbf{x}\in\bigcap_{i=1}^k\psi_i(Q_i)
=\psi(Q),
\]
and thus $\psi(Q)$ is a non-empty cell of $\mathscr{T}^h$. This proves that the map
\[
Q\longmapsto\psi(Q)
\]
is well defined.

Conversely, let
\[
\tilde Q=\bigcap_{i=1}^k\tilde Q_i
\]
be a non-empty cell of $\mathscr{T}^h$, where each $\tilde Q_i$ is a non-empty cell of $\mathcal{T}_i^h$. Choose any $\mathbf{x}\in\tilde Q$ in the relative interior of each $\tilde Q_i$. For each $i\in[k]$, the cell $\tilde Q_i$ corresponds to the face~$Q_i$ of $P_i$ minimizing
\[
-h(\{e_i\}\times R)+\langle R,\mathbf{x}\rangle,
\qquad
R\in\operatorname{Ver}(P_i).
\]
The minimum values need not coincide for different $i$. However, choosing and adding constants $y_1,\dots,y_k$ so that these minima become equal, and setting $\mathbf{y}=(y_1,\dots,y_k)$, the upper face minimizing $(-1,\mathbf{y},\mathbf{x})$ projects precisely to
\[
Q=
\operatorname{conv}\bigl((\{e_1\}\times Q_1)\cup\cdots\cup(\{e_k\}\times Q_k)\bigr).
\]
Consequently,
\[
\psi(Q)
=\bigcap_{i=1}^k\psi_i(Q_i)
=\bigcap_{i=1}^k\tilde Q_i
=\tilde Q.
\]
Therefore every non-empty cell of $\mathscr{T}^h$ is the image of a face of $\mathcal{S}^h$ containing at least one vertex from each copy $\{e_i\}\times P_i$. Thus, the map $\psi$ is surjective.

To prove injectivity, observe that, for a fixed $\mathbf{x}$, the vector $\mathbf{y}$ is uniquely determined up to adding a multiple of $(1,\dots,1)$. Adding the same constant to all coordinates of $\mathbf{y}$ adds the same constant to every value of the linear functional \eqref{eq_minimizedfaces}, and therefore does not change its minimizing face. Hence the face $Q$ is uniquely determined by the cell $\tilde Q=\psi(Q)$, showing that $\psi$ is injective.

Finally, suppose
\[
Q=
\operatorname{conv}\bigl((\{e_1\}\times Q_1)\cup\cdots\cup(\{e_k\}\times Q_k)\bigr)
\subseteq
Q'=
\operatorname{conv}\bigl((\{e_1\}\times Q_1')\cup\cdots\cup(\{e_k\}\times Q_k')\bigr).
\]
Then $Q_i\subseteq Q_i'$ for every $i\in[k]$. Hence every point of $\psi_i(Q_i')$ satisfies the defining inequalities for $\psi_i(Q_i)$, and therefore
\[
\psi_i(Q_i')\subseteq\psi_i(Q_i).
\]
Intersecting over all $i$ gives
\[
\psi(Q')
=
\bigcap_{i=1}^k\psi_i(Q_i')
\subseteq
\bigcap_{i=1}^k\psi_i(Q_i)
=
\psi(Q),
\]
showing that $\psi$ is order reversing.
\end{proof}

\end{document}
